\documentclass[10pt,a4paper]{amsart}
\usepackage[margin=19mm]{geometry}
\usepackage{amsmath,amssymb,mathtools}
\usepackage[scr=rsfs,cal=euler]{mathalfa}
\usepackage{xfrac}
\usepackage[unicode,hidelinks,bookmarks=false]{hyperref}
\newcommand{\ie}{i.e.\ }
\newcommand{\cf}{cf.\ }
\newcommand{\resp}{respectively\ }
\newcommand{\Subsection}{\subsection}
\providecommand{\nobreakdash}{\nobreak\hbox{-}}
\setlength{\emergencystretch}{2em}
\multlinegap0pt
\usepackage{amssymb}
\usepackage{enumerate}
\usepackage{xcolor}
\usepackage{bm}
\newtheorem{condition}{Condition}
\newtheorem{theorem}{Theorem}
\newtheorem{corollary}{Corollary}
\newcommand\CG{{\mathcal G}}
\newcommand\CR{{\mathcal R}}
\newcommand\E{{\mathbb{E}}}
\renewcommand\P{{\mathbb{P}}}
\newcommand\R{{\mathbb{R}}}
\newcommand\Z{{\mathbb{Z}}}
\newcommand\eps{{\varepsilon}}
\title{The anti-concentration phenomenon with respect to random permutations}
\author{Viet H. Do, Hoi H. Nguyen, Kiet H. Phan, Tuan Tran, Van H. Vu}
\date{Source: arXiv:2512.21779v2 (CC BY 4.0)}
\begin{document}
\maketitle

\noindent\emph{Source and licence.} The anti-concentration phenomenon with respect to random permutations. Version arXiv:2512.21779v2 (CC BY 4.0), distributed by arXiv under the Creative Commons Attribution 4.0 International licence, http://creativecommons.org/licenses/by/4.0/, as recorded in the arXiv licence metadata for this exact version and shown on the abstract page https://arxiv.org/abs/2512.21779v2. Full licence notice: \texttt{licenses/arxiv-2512.21779-CC-BY-4.0.txt}.\par\smallskip

\noindent\textit{Editorial scope.} Counted unit: the article's theorem thm:cont:3/2 (source0.tex lines 822--844). The article's complete original proof of it is delivered with it (source0.tex lines 1952--2089, one \texttt{proof} environment of 138 lines, which the article places away from the statement). No mathematics is written, repaired or re-derived: the statement and the proof are the article's own lines, in the article's own notation, and no retained line is substituted or re-flowed.  Retained uncounted as the source-local context the proof needs: lines 771--783; lines 798--801; lines 1313--1319; lines 1376--1379; lines 1684--1687.  Nothing was omitted from any retained span, so the package carries no omission record.  No retained line contains a citation, so the wrapper carries no bibliography and no bibliography entry.  No retained line contains a figure environment, an image-inclusion command or drawing code, so no image is delivered and the package declares no asset dependency.  Editorial formatting only, in addition: an AMS \texttt{amsart} preamble that restates the article's own declarations and macros used by the retained lines, namely the environments \texttt{condition}, \texttt{theorem}, \texttt{corollary} on separate counters, together with the standard packages those lines need.  The retained environments print in this document as Condition 1, Theorem 1, Corollary 1 in order of appearance, whereas the article prints the same items with the numbering of its own full document.  The complete native source of the article as distributed in the arXiv e-print for this exact version is bundled unchanged as \texttt{sources/191686/source0.tex}.
\begin{condition}[Non-degeneracy]\label{cond:separation}
We say that a sequence $w_1,\dots,w_n$ satisfying
\[
\sum_{i=1}^n w_i = 0
\quad\text{and}\quad
\sum_{i=1}^n w_i^2 = 1
\]
is \emph{non-degenerate} if
\begin{equation}\label{eqn:separation}
|w_i-w_j| \;\le\; \frac{1}{A\sqrt{\log n}}
\end{equation}
for all distinct $i,j$ and for some constant $A$.
\end{condition}

\begin{equation}\label{eqn:1:u}
\P\Bigl(\bigl|S_{\pi}-L\bigr|\le \frac{1}{n}\Bigr)
= O_{A}\Big(\frac{1}{n}\Big).
\end{equation}

\begin{corollary}\label{cor:Roos}
Assume that $S_{\pi} = \sum_{k=1}^{n} a_{k \pi(k)}$.
Then its characteristic function $\varphi(t) = \E e^{it S_{\pi}}$ satisfies
\[
|\varphi(2\pi t)| \le \exp\Big(-\frac{1}{2n^{3}} \sum_{1\le i,j,k,l\le n} \|t (a_{ik} - a_{jk} -a_{il} + a_{jl})\|^2_{\R/\Z}\Big).
\] 
\end{corollary}

\begin{corollary}\label{cor:wraparound} 
Let $0 < \delta < 1$ and $I \subset \{-n,\dots,n\}$ with $|I| \geq \delta n$. For every constant $C > 0$ sufficiently large in terms of $\delta$, and for all $\frac{1}{Cn} \leq |b| \leq \frac{1}{C}$ and for any $b_0 \in \R$ we have
$$\sum_{r \in I} \| b r +b_0 \|_{\R/\Z}^2 =\Theta_{C} (n).$$
\end{corollary}

\begin{equation}\label{eqn:int:1}
\P(|X| \le 1) =O\Big(\int_{|t| \le 1} |\varphi_{S}(t)| \, dt\Big)
=O\Big(\int_{|t| \le 1} \exp\Big\{ -\frac{c_{\delta}}{n^{2}} \sum_{\substack{1 \le i,j \le n \\ r\in \CR}} \| t (w_{i} - w_{j}) r \|_{\R/\Z}^{2} \Big\} dt\Big).
\end{equation}

\begin{theorem}[New treatment at scale $n^{-3/2}$]\label{thm:cont:3/2}
Let $0<\varepsilon<1/2$ and $\delta>0$ be given. Suppose that the sequence
$(w_1,\ldots,w_n)$ satisfies Condition~\ref{cond:separation} for some
sufficiently large constant $A>0$. Furthermore, assume that no interval of
length $\varepsilon/\sqrt{n}$ contains more than $(1-\varepsilon)n$ of the
values $w_i$. Let $I\subset[n]$ be any subset with $|I|\ge \delta n$, and
consider a sequence $(v_1,\ldots,v_n)$ that is partially specified by
\[
v_i = i/n \quad \text{for all } i\in I.
\]
Then, for every $L\in\R$, we have the uniform bound
\begin{equation}\label{eqn:3/2:u}
\P\Bigl(\bigl|S_{\pi}-L\bigr|\le \frac{1}{n^{3/2}}\Bigr)
= O_{A}\Big(\frac{1}{n^{3/2}}\Big).
\end{equation}
If, in addition, there exists a constant $\widetilde{B}>0$ such that
$|v_i|\le \widetilde{B}$ for all $i\in[n]$, then we obtain the sub-gaussian bound
\begin{equation}\label{eqn:3/2:L}
\P\Bigl(\bigl|S_{\pi}-L\bigr|\le \frac{1}{n^{3/2-\eps}}\Bigr)
= O\Big(\frac{1}{n^{3/2-\eps}} e^{-\Theta(L^{2})}\Big).
\end{equation}
Here, the implied constants depend only on $A$ and $\widetilde{B}$.
\end{theorem}

\begin{proof}[Proof of \eqref{eqn:3/2:u} of Theorem \ref{thm:cont:3/2}]
Let $\Delta>0$ be a constant chosen sufficiently large in terms of $\eps$ and $\delta$.\footnote{Only the small-$|t|$ regime considered below requires $\Delta$ to be a constant. The other regimes hold for \emph{any} $\Delta = \Delta(n)$ sufficiently large with respect to $\eps$ and $\delta$.} Define
\[
S' = \sum_{i} \frac{n^{3/2}}{\Delta}w_{i} v_{\pi(i)} =\frac{n^{1/2}}{\Delta} S
\quad\text{and}\quad
X' = S' - \frac{n^{3/2}}{\Delta}L.
\]
To prove \eqref{eqn:3/2:u}, it suffices to show $\P(|X'| \le 1)=O_{\eps,\delta,\Delta,A}(\frac{1}{n^{3/2}})$.
Using Esseen's estimate together with Corollary \ref{cor:Roos}, we can write
\[
\sup_{x} \P(|S' - x| \le 1)
= O\Big(\int_{|t| \le 1} |\varphi_{S'}(t)| \, dt\Big)
=O\Big(\int_{|t| \le 1} \exp\Big\{- \frac{1}{2n^{3}} \sum_{i,j,k,l} \Big\|\frac{t n^{3/2}}{\Delta} (w_{i}-w_{j})(v_{k}-v_{l})\Big\|_{\R/\Z}^{2}\Big\} \, dt\Big).
\]
Recalling that $v_i = i/n$ for $i\in I$, and following the argument from the proof of \eqref{eqn:1:u}, the exponent on the right-hand side can be bounded from below by
\[
\frac{c_{\delta}}{n^{2}} \sum_{\substack{1 \le i,j \le n \\ r\in \CR}}
\Big\| \frac{t n^{1/2}}{\Delta} (w_{i} - w_{j}) r \Big\|_{\R/\Z}^{2},
\]
and so
\begin{equation}\label{eqn:int:3/2}
\P(|X'| \le 1)  
= O\Big(\int_{|t| \le 1} \exp\Big\{ -\frac{c_{\delta}}{n^{2}} \sum_{\substack{1 \le i,j \le n \\ r\in \CR}} \Big\| \frac{t n^{1/2}}{\Delta}(w_{i} - w_{j}) r \Big\|_{\R/\Z}^{2} \Big\} dt\Big).
\end{equation}

We note that this differs from \eqref{eqn:int:1} in that we have an extra factor of $n^{1/2}$ in the exponent.  
As such, our case analysis for $t$ will be different, and we will need more information in addition to \eqref{eqn:separation} about the $w_{i}$ (as assumed in Theorem \ref{thm:cont:3/2}).

We discard pairs $(i,j)$ for which $|w_{i}-w_{j}|$ is much smaller than $1/\sqrt{n}$, and set
\[
\CG = \{ (i,j) : |w_{i}-w_{j}| \ge \eps / 2\sqrt{n} \}.
\]
Since there is no interval of length $\eps/\sqrt{n}$ containing at least $(1-\eps)n$ elements from $w_{1},\dots,w_{n}$, the number of pairs $(i,j)$ such that $|w_{i}-w_{j}| < \eps/2\sqrt{n}$ is at most $(1-\eps)n^{2}$.
Thus, we have
\begin{equation}\label{eqn:3/2:u:starting}
|\CG|\ge \eps n^2, \qquad \sum_{(i,j)\in \CG} (w_{i}-w_{j})^{2} =\Theta_{\eps}(n).
\end{equation}

\underline{\bf Intermediate $|t|$.}  
We first focus on the range
\[
\frac{(\sqrt{A\log n})\Delta}{n^{3/2}} \le |t| \le \frac{\Delta}{n^{1/2}}.
\]

Now, for $\log \big(\frac{2}{\eps}\big)\le k \le \log\Big( \frac{\sqrt{n}}{A \sqrt{\log n}} \Big)+1$ (so that $\frac{2^{k-1}}{\sqrt{n}} \le \frac{1}{A\sqrt{\log n}}$ and $\frac{2^{k}}{\sqrt{n}} \ge \frac{\eps}{2\sqrt{n}}$), let $\CG_{k}$ be the collection of pairs $(i,j)$ for which
\[
D_{k-1} = \frac{2^{k-1}}{\sqrt{n}} < |w_{i}-w_{j}| \le \frac{2^{k}}{\sqrt{n}} = D_{k}.
\]
Then we see that $\CG=\bigcup_{k}\CG_{k}$ and
\[
\sum_{k} D_{k}^{2} |\CG_{k}| =\Theta_{\eps}(n).
\]

Given a fixed pair $(i,j)\in \CG_{k}$, we consider the sum $\sum_{r\in \CR} \|\frac{t n^{1/2}}{\Delta} (w_{i}-w_{j})r\|_{\R/\Z}^{2}$. For $\frac{t n^{1/2}}{\Delta}D_{k}\le \frac{1}{C n}$, we have $|\frac{t n^{1/2}}{\Delta} (w_{i}-w_{j})r|\le \frac{1}{C}<1$, so
\[
\sum_{r\in \CR}\Big\|\frac{t n^{1/2}}{\Delta} (w_{i}-w_{j})r\Big\|^{2}_{\R/\Z} = \sum_{r\in \CR}\Big|\frac{t n^{1/2}}{\Delta} (w_{i}-w_{j})r\Big|^2 =\Theta_{\delta} \Big(\frac{t^{2}n^{4}D^{2}_{k}}{\Delta^{2}}\Big).
\]
On the other hand, if $\frac{tn^{1/2}}{\Delta}D_{k}> \frac{1}{Cn}$, then $\frac{1}{2Cn}\le |\frac{tn^{1/2}}{\Delta}(w_{i}-w_{j})|\le \frac{1}{A\sqrt{\log n}}$, and by Corollary \ref{cor:wraparound}, we get
\[
\sum_{r\in \CR} \Big\|\frac{tn^{1/2}}{\Delta} (w_{i}-w_{j})r\Big\|_{\R/\Z}^{2} =\Theta_{C}(n).
\]
From the above discussion, we conclude that
\[
\sum_{i,j,r} \Big\|\frac{tn^{1/2}}{\Delta} (w_{i}-w_{j})r\Big\|_{\R/\Z}^{2}
=\Omega_{\delta,C}(1)\cdot
\Big(
\sum_{k: \, \frac{|t|n^{1/2}}{\Delta}D_{k} \le \frac{1}{Cn}} \frac{t^2 n^{4} D_{k}^{2}}{\Delta^2} |\CG_{k}|
+ \sum_{k: \, \frac{|t|n^{1/2}}{\Delta}D_{k}>\frac{1}{Cn}} n |\CG_{k}|\Big).
\]

To lower bound the right-hand side, we distinguish two cases.

\medskip
\noindent\textbf{Case 1:} $\sum_{k: \, \frac{|t|n^{1/2}}{\Delta}D_{k} \le \frac{1}{Cn}} D_{k}^{2} |\CG_{k}| =\Theta_{\eps}(n)$.

Since $t^2 \ge \frac{(A\log n)\Delta^{2}}{n^{3}}$, we have
\[
\sum_{k: \, \frac{|t|n^{1/2}}{\Delta}D_{k} \le \frac{1}{Cn}} \frac{t^2 n^{4} D_{k}^{2}}{\Delta^{2}} |\CG_{k}|
=\Omega_{\eps}(An^{2}\log n).
\]

\medskip
\noindent\textbf{Case 2:} $\sum_{k:\, \frac{|t|n^{1/2}}{\Delta}D_{k} > \frac{1}{Cn}} D_{k}^{2} |\CG_{k}| =\Theta_{\eps}(n)$.

As $D^{2}_{k}\le 1/(A^2\log n)$, we obtain $\sum_{k:\, \frac{|t|n^{1/2}}{\Delta}D_{k} > \frac{1}{Cn}} |\CG_{k}|= \Omega_{\eps}(A^2 n \log n)$,
which implies
\[
\sum_{k:\, \frac{|t|n^{1/2}}{\Delta}D_{k} > \frac{1}{Cn}} n |\CG_{k}|
=\Omega_{\eps}(A^2 n^{2}\log n).
\]

Therefore, in both cases,
\[
\sum_{i,j,r} \|t n^{1/2} (w_{i}-w_{j})r\|_{\R/\Z}^{2} =\Omega_{\eps,\delta,C}(A n^{2} \log n).
\]
Assuming $A$ is sufficient large relative to $\eps,\delta$, and $C$, this implies that for $\frac{ \sqrt{A\log n}}{n^{3/2}} \le |t| \le \frac{1}{n^{1/2}}$, we have
\[
\frac{c_{\delta}}{n^{2}}\sum_{i,j,r} \| t (w_{i} - w_{j}) r \|_{\R/\Z}^{2} \ge 2\sqrt{A} \log n, \qquad |\varphi_{S'}(t)| \le n^{-2\sqrt{A}}.
\]


\underline{\bf Large $|t|$.} We now focus on the range 
$$\frac{\Delta}{n^{1/2}} \le |t|\le 1.$$  

Let $\CG_{0}$ be the set obtained from $\CG$ by removing pairs $(i,j)$ with $|w_{i} - w_{j}| > 2/\sqrt{\eps n}$. Since $\sum_{i,j}(w_{i}-w_{j})^2=2n$, at most $(\eps/2)n^2$ pairs were removed. From \eqref{eqn:3/2:u:starting}, we see that $|\CG_{0}|\ge |\CG|-(\eps/2)n^2\ge (\eps/2)n^2$. For every $(i,j)\in \CG_{0}$, as $\frac{\eps}{2\sqrt{n}} \le |w_{i} - w_{j}| \le \frac{2}{\sqrt{\eps n}}$, we have
\[
\frac{\eps}{2\sqrt{n}}\le \Big|\frac{t n^{1/2}}{\Delta} (w_{i}-w_{j})\Big|\le \frac{2}{\Delta \sqrt{\eps}}.
\]
Assuming $\Delta$ is sufficiently large in terms of $\eps$ and $\delta$, Corollary \ref{cor:wraparound} then gives
\[
\frac{c_{\delta}}{n^{2}}
\sum_{\substack{(i,j)\in \CG_{0}\\ r\in \CR}}
\Big\| \frac{t n^{1/2}}{\Delta} (w_i-w_j) r \Big\|_{\R/\Z}^{2}
= \Theta_{\delta}\!\Big(\frac{|\CG_{0}|}{n}\Big)
= \Theta_{\delta,\varepsilon}(n).
\]

\underline{\bf Small $|t|$.}  
It remains to consider 
$$|t| \le \frac{(\sqrt{A\log n})\Delta}{n^{3/2}}.$$

In this case, as $|w_{i}- w_{j}| \le 1/(A \sqrt{\log n})$, we have
\[
\Big|\frac{t n^{1/2}}{\Delta} (w_{i}-w_{j})r\Big|\le \frac{1}{\sqrt{A}}<1.
\]
Therefore,
\[
\frac{c_{\delta}}{n^{2}} \sum_{i,j,r} \Big\|\frac{t n^{1/2}}{\Delta} (w_{i}-w_{j})r\Big\|_{\R/\Z}^{2}
= \frac{c_{\delta}}{n^{2}} \sum_{i,j,r} \Big|\frac{t n^{1/2}}{\Delta}(w_{i}-w_{j})r\Big|^{2}
=\Theta_{\delta,\Delta}(t^{2} n^{3}).
\]
Hence,
\[
\int_{|t|\le \frac{(\sqrt{A\log n})\Delta}{n^{3/2}}} \exp\Big\{- \frac{c_{\delta}}{n^{2}} \sum_{\substack{1\le i,j \le n \\r \in \CR}} \Big\| \frac{t n^{1/2}}{\Delta} (w_{i}-w_{j}) r\Big\|_{\R/\Z}^{2}\Big\}\Big)
\le \int_{\R} e^{-\Theta(t^{2} n^{3})}\, dt
= O(n^{-3/2}).
\]
\end{proof}

\end{document}
